\chapter*{Abstract\markboth{Abstract}}
\addcontentsline{toc}{chapter}{Abstract}

Fetal diffusion magnetic resonance imaging (dMRI) is highly sensitive to fetal motion and other acquisition-related artifacts. These artifacts can reduce image quality, affect later reconstruction and analysis, and may make parts of the acquired data unusable. Reliable quality control is therefore important for identifying corrupted image data and improving the overall usability of fetal dMRI.

In this thesis, an artificial intelligence-based quality-control framework was developed for fetal dMRI. The workflow is based on automatic localization of the fetal brain followed by artifact assessment. Two different processing strategies were investigated: a volumetric 3D approach and a lighter 2D slice-wise approach. The 3D pipeline evaluates the complete MRI volume, while the 2D pipeline processes individual slices and combines their predictions to obtain a final quality decision for the complete dynamic. The study also evaluates brain segmentation, artifact classification, slice-to-dynamic aggregation, and computational performance.

The experiments showed that the 3D approach provided more reliable artifact detection because it uses information from the complete volume. The 2D approach reduced the computational effort and offered a faster processing pathway, although its performance depended strongly on how the slice-level predictions were combined. Further testing of aggregation strategies improved the sensitivity of the 2D approach while maintaining good specificity.

Overall, this study demonstrates the potential of AI-based quality control for fetal dMRI and shows how volumetric and slice-wise processing can be used for automated artifact assessment. The comparison highlights the trade-off between reliable artifact detection and computational efficiency and identifies important factors for future quality-control systems, including robust fetal-brain localization, suitable aggregation of slice-level predictions, and reliable quality labels.
